\section{Using the conclusions and locating their limits}\label{sec:use}
The preceding proofs provide both a bound to apply and intermediate statements to reuse. This section distinguishes those uses from numerical evidence or an assertion of optimal constants.

\subsection{An explicit route from information to a coordinate}
Throughout this section, $n\ge1$. Let
\[
 \mathcal D_p(f)=\ln2-h(p)-\II(f(X);Y).
\]
For $0<p<1/2$, the proof gives the stronger piecewise coefficient
\begin{equation}\label{eq:usable-coefficient}
 c_p=\begin{cases}
 \dfrac{9p(1-2p)^2}{10000},&0<p<1/20,\\[4pt]
 \dfrac{(1-2p)^2}{2000},&1/20\le p<1/2,
 \end{cases}
 \qquad \mathcal D_p(f)\ge c_p\,\delta(f).
\end{equation}
These are precisely~\eqref{eq:low-stability} and~\eqref{eq:complement-stability}; at the shared endpoint we use the larger coefficient. The change of formula reflects the proof, not a discontinuity of mutual information.

\begin{corollary}[Recovering a coordinate from a near-optimal value]\label{cor:recovery}
If $\mathcal D_p(f)\le\epsilon$ at fixed $0<p<1/2$, there are $i$ and $\sigma\in\{-1,1\}$ such that
\begin{equation}\label{eq:recovery}
 \Pp\{2f(X)-1=\sigma X_i\}\ge
 1-\min\{1/2,\epsilon/c_p\}.
\end{equation}
For $\epsilon\ge0$, the coordinate can be chosen to maximize $|\E[(2f-1)X_i]|$, with its sign chosen to agree with that coefficient.
\end{corollary}
\begin{proof}
For $F=2f-1$,
$\Pp(F\ne\sigma X_i)=(1-\sigma\E(FX_i))/2$.
Minimizing this expression identifies the stated coordinate and gives
$\delta(f)=(1-\max_i|\E(FX_i)|)/2$. Now use~\eqref{eq:usable-coefficient} and $\delta(f)\le1/2$.
\end{proof}

For example, at $p=1/4$ one has $c_p=1/8000$. An information gap of at most $10^{-6}$ nats therefore guarantees disagreement on at most $1/125$ of the inputs, in every finite dimension. The coefficient is deliberately conservative. The corollary becomes nontrivial only when a valid upper bound on the information gap is sufficiently small. It neither supplies a statistical estimator of that gap nor claims sharp numerical performance or an efficient learning algorithm.

There is also an elementary converse, requiring no new CK machinery.
\begin{proposition}[Closeness implies a small information gap]\label{prop:continuity}
For $0<p<1/2$,
\begin{equation}\label{eq:stability-sandwich}
 c_p\delta(f)\le\mathcal D_p(f)\le h(\delta(f)).
\end{equation}
\end{proposition}
\begin{proof}
The lower bound is~\eqref{eq:usable-coefficient}. Choose a nearest coordinate indicator $g$ and put $E=f\mathbin{\oplus}g$, the error bit. The chain rule gives
\[
 \II(g;Y)-\II(f;Y)
 =\II(g;Y\mid f)-\II(f;Y\mid g)
 \le\HH(g\mid f)=\HH(E\mid f)\le h(\delta(f)).
\]
Since $\II(g;Y)=\ln2-h(p)$, this is the upper bound.
\end{proof}
The two bounds give a qualitative equivalence between near optimality and nearness to a coordinate at fixed nontrivial noise. Their rates do not match: the lower bound is linear, whereas $h(\delta)$ is of order $\delta\ln(1/\delta)$ near zero. No sharp stability rate is asserted.

\subsection{Three examples that explain the hypotheses}\label{sec:examples}
\paragraph{Constants and coordinates.}
A constant output has $V(\mu)=0$, zero conditional entropy, and zero information. It saturates the variance profile but not the information optimum when $0<p<1/2$. A coordinate has $V=1$ and conditional entropy $h(p)$, and saturates both bounds. This distinction explains the separate output-bias term in Corollary~\ref{cor:low}.

\paragraph{Parity: zero mean gap does not mean zero posterior gap.}
For $F(x)=\prod_{i=1}^k x_i$, the noisy posterior mean is $T_\rho F=\rho^kF$. Thus
\begin{equation}\label{eq:parity-example}
 \HH(f\mid Y)=\eta(\rho^k),\qquad
 \Delta_p(f)=\mathcal D_p(f)=\eta(\rho^k)-\eta(\rho).
\end{equation}
For $k\ge2$, all mass is at degree $k$, and $\delta(f)=1/2$. The two-bit case is especially useful: both sections are balanced, so $d=0$, but their posterior difference has constant magnitude $\rho$. In the exact deficit identity~\eqref{eq:deficit}, the two child deficits, the triangle term, and the bias surplus vanish. All remaining entropy comes from the scalar mixing slack.

\paragraph{AND: unequal section means are unavoidable.}
For $f(x)=\one\{x_1=x_2=1\}$,
\[
 F(x)=\frac{-1+x_1+x_2+x_1x_2}{2},\qquad
 \mu=\frac14,\quad m^2=\frac14,\quad W_1(F)=\frac12,\quad
 \sum_{k\ge2}W_k(F)=\frac14,\quad\delta(f)=\frac14.
\]
The conditional entropy can be calculated without Fourier theory:
\begin{equation}\label{eq:and-example}
 \HH(f\mid Y)=\frac14\bigl[h(p^2)+2h(p(1-p))+h((1-p)^2)\bigr].
\end{equation}
The four terms correspond to the four equally likely noisy observations. Fixing one coordinate gives a constant section and a coordinate section, not two balanced functions. The example checks both the meaning of the variance loss and the factor-four conversion between indicator and sign Fourier weights.

\subsection{Why the uniform scalar margin has the scale it does}
\begin{proposition}[An upper limit on a uniform quadratic surplus]\label{prop:scalar-barrier}
Fix $0<p<1/2$. If a number $\gamma(p)$ satisfies
$S_{p,d}(a,b)\ge\gamma(p)(a-b)^2$ for all $a,b\in[0,1]$ and independently all $d\in[0,1/2]$, then
\begin{equation}\label{eq:scalar-barrier}
 \gamma(p)\le2p(1-p).
\end{equation}
\end{proposition}
\begin{proof}
Take $d=0$, $a=(1-t)/2$, and $b=(1+t)/2$. The slack is then $\eta(\rho t)-\eta(t)$. Since $\eta'(0)=0$ and $\eta''(0)=-1$, Taylor's formula gives
\[
 \lim_{t\downarrow0}\frac{\eta(\rho t)-\eta(t)}{t^2}
 =\frac{1-\rho^2}{2}=2p(1-p).
\]
Apply the hypothesized bound and take this limit.
\end{proof}
For $0<p\le1/20$, the best possible coefficient in this particular all-posterior statement therefore lies between $p/500$ and $2p(1-p)$. In particular it cannot be bounded below by $Cp\ln(1/p)$ for one positive constant $C$ and all sufficiently small $p$.

This does not preclude a logarithm for actual posteriors attached to a fixed Boolean function. For two-bit parity,~\eqref{eq:parity-example} reads
\[
 \Delta_p(f)=h(2p(1-p))-h(p)
 =p\ln(1/p)+(1-2\ln2)p+O\bigl(p^2\ln(1/p)\bigr).
\]
Here the posteriors approach deterministic endpoints as $p\downarrow0$. The uniform scalar coefficient and the deficit along such a family of posteriors are different objects. This fixed two-bit expansion is not a uniform asymptotic theorem over dimension.

\subsection{What can be taken from the proof}
Proposition~\ref{prop:induction} can be used whenever a two-posterior inequality supplies a profile coefficient $c$. Lemma~\ref{lem:product} supplies an entropy-product comparison without reference to a Boolean dimension. Proposition~\ref{prop:certificate} separates a spectral cap from the entropy envelope needed to convert it into information. The induction behind~\eqref{eq:generic-spectral} explains how a local surplus pays for precisely the spectral mass lost under restriction.

These statements suggest concrete questions rather than a blanket claim of generality. How large can the uniform scalar surplus be within the upper limit~\eqref{eq:scalar-barrier}? Can a better cap or a better envelope improve the constants in~\eqref{eq:generic-certificate}? What additional statistics of realizable posteriors recover information discarded by the all-pair scalar bound? None of the present results covers biased product inputs or asymmetric channels: equal section weights, uniform output of the channel, and the sign-cube Fourier basis enter explicitly.
